\section*{Chapter \ref{ch:b1:comparative-statics} --- Comparative Statics}
% Numbers recomputed in code/bachelor-1/03-comparative-statics/exercises/solutions.py (tested).

\begin{solution}{exo:b1:comparative-statics:ce}
\begin{enumerate}
  \item $\ln q = \ln 50 - 0.5\ln p$, so $\varepsilon = -0.5$ at every price.
  \item $q = 50/\sqrt 4 = 25$. A $10\%$ rise predicts $-0.5 \times 10\% = -5\%$,
    so $q \approx 23.75$. Exactly, $50/\sqrt{4.4} = 23.84$: the first-order
    prediction is close for a $10\%$ change.
\end{enumerate}
\end{solution}

\begin{solution}{exo:b1:comparative-statics:read-slope}
\begin{enumerate}
  \item $2^{1.25} = 2.38$: the richer country uses about 2.4 times as much energy per person.
  \item $\ln(292.7/21.85)/\ln(69121/7702) = \ln 13.40/\ln 8.974 = 2.595/2.194 = 1.18$,
    a little below the slope of the line through all 183 countries.
\end{enumerate}
\end{solution}

\begin{solution}{exo:b1:comparative-statics:demand-shift}
$F(p, a) = (a - bp + u^d) - (c + dp + u^s)$ has $F_p = -(b + d)$ and $F_a = 1$, so
$\dd p^*/\dd a = 1/(b + d) > 0$. Then $q^* = c + dp^* + u^s$ gives
$\dd q^*/\dd a = d/(b + d) > 0$: a rise in demand raises both, as in
\cref{prop:b1:models-and-data:linear}.
\end{solution}

\begin{solution}{exo:b1:comparative-statics:spending}
By \cref{prop:b1:comparative-statics:rules},
$\varepsilon_{pq, p} = \varepsilon_{p,p} + \varepsilon_{q,p} = 1 + \varepsilon_d$.
With $\varepsilon_d = -0.3$, a $10\%$ rise in price raises spending by about
$0.7 \times 10\% = 7\%$.
\end{solution}

\begin{solution}{exo:b1:comparative-statics:implicit-tax}
\begin{enumerate}
  \item $F(p, \tau) = 100 - 2p - 3(p - \tau)$, with $F_p = -5$ and $F_\tau = 3$.
  \item $\dd p/\dd\tau = 3/5 = 0.6$; $\dd(p - \tau)/\dd\tau = -0.4$;
    $\dd q/\dd\tau = D' \times 0.6 = -2 \times 0.6 = -1.2$.
  \item $5p = 100 + 3\tau$ gives $p = 20 + 0.6\tau$ and $q = 100 - 2p = 60 - 1.2\tau$.
\end{enumerate}
\end{solution}

\begin{solution}{exo:b1:comparative-statics:fuel}
\begin{enumerate}
  \item $1.2/(1.2 + 0.3) = 0.8$.
  \item The pump price rises by $0.8 \times 10 = 8$ cents; sellers bear the
    other 2 cents, a fifth of the tax.
  \item As $\varepsilon_s$ grows, $\varepsilon_s/(\varepsilon_s - \varepsilon_d)$ tends
    to 1: buyers bear the whole tax.
\end{enumerate}
\end{solution}

\begin{solution}{exo:b1:comparative-statics:envelope-simple}
$f(x, a) = ax - x^2$ has $f_x = a - 2x = 0$ at $x^* = a/2$, and
$V(a) = a^2/2 - a^2/4 = a^2/4$. So $V'(a) = a/2$, and
$\partial f/\partial a = x$, which is $a/2$ at $x^*$: the theorem holds.
\end{solution}

\begin{solution}{exo:b1:comparative-statics:two-years}
\begin{enumerate}
  \item In both years energy use rises more than in proportion to income
    across countries (slope above 1), but less so in 2019: a $1\%$ income gap
    went with $1.25\%$ more energy use, against $1.35\%$ in 1990. The
    cross-section flattened.
  \item So that the change in slope is not caused by countries entering the
    data (the 31 countries with 2019 data but none for 1990), only by the
    same countries' changes.
  \item No. It describes two cross-sections. The response of one country's
    energy use to its own growth would need its own changes over time,
    compared with other countries' changes, and a reason why those income
    changes were not themselves caused by energy.
\end{enumerate}
\end{solution}

\begin{solution}{exo:b1:comparative-statics:which-curve}
Neither, as it stands. It uses differences between countries in one year.
For it to measure a country's response of energy demand to income, it would
need (i) no other difference between rich and poor countries that changes
energy use --- climate, the price of energy, the share of heavy industry ---
to be correlated with income, and (ii) income not to be caused by energy:
countries rich because they sell oil and gas sit high on both axes for that
reason.
\end{solution}

\begin{solution}{exo:b1:comparative-statics:envelope-proof}
Differentiate $V(\theta) = f(x^*, y^*, \theta)$:
$V' = f_x x^{*\prime} + f_y y^{*\prime} + f_\theta$. The constraint holds for
every $\theta$, so $g_x x^{*\prime} + g_y y^{*\prime} + g_\theta = 0$. The \omterm{thm:b1:constrained-optimisation:lagrange}{Lagrange
conditions} are $f_x = \lambda^* g_x$ and $f_y = \lambda^* g_y$, so
$f_x x^{*\prime} + f_y y^{*\prime} = \lambda^*(g_x x^{*\prime} + g_y y^{*\prime})
= -\lambda^* g_\theta$. Hence $V' = f_\theta - \lambda^* g_\theta$, which is
$\partial\mathcal L/\partial\theta$ for $\mathcal L = f - \lambda g$.
\end{solution}

\begin{solution}{exo:b1:comparative-statics:monopoly}
\begin{enumerate}
  \item Profit is $(p - c)Ap^{-\varepsilon}$. Its derivative,
    $Ap^{-\varepsilon} - \varepsilon(p - c)Ap^{-\varepsilon - 1}$, is zero when
    $p = \varepsilon(p - c)$, that is $p = c\,\varepsilon/(\varepsilon - 1)$.
  \item With the tax the seller's cost per unit is $c + \tau$, so
    $p = (c + \tau)\varepsilon/(\varepsilon - 1)$ and
    $\dd p/\dd\tau = \varepsilon/(\varepsilon - 1)$, which exceeds 1 for every
    $\varepsilon > 1$.
  \item The markup factor is $2/(2 - 1) = 2$: $p = 20$ before and $22$ after, a
    rise of 2 for a tax of 1.
\end{enumerate}
\end{solution}

\begin{solution}{exo:b1:comparative-statics:refute}
\begin{enumerate}
  \item A refutation: in a market with many sellers, with costs and demand
    otherwise unchanged (other markets nearby, untaxed, as a comparison),
    the price rises by more than the tax, or falls. Not a refutation: the
    price rises by more than the tax in the same month as a rise in crude oil
    prices, or in a market with one or two sellers, where the model's
    assumptions do not hold.
  \item It contradicts the assumption that buyers react to the tax as to any
    part of the price. For a tax included in the posted price, buyers see it,
    and the prediction of pass-through is unchanged; the finding warns that
    the same tax, added at the till, would be passed on and reacted to
    differently.
\end{enumerate}
\end{solution}

\begin{solution}{pb:b1:comparative-statics:cement}
\textbf{1.} $D(100) = 1000 \times 1^{-0.5} = 1000$ and $S(100) = 1000 \times 1^{1.2} = 1000$.

\textbf{2.} $F(p, \tau) = 1000(p/100)^{-0.5} - 1000\bigl((p - \tau)/100\bigr)^{1.2} = 0$.

\textbf{3.} $D'(p) = -0.5 \times 10 \times (p/100)^{-1.5}$, so $D'(100) = -5$;
$S'(p_s) = 1.2 \times 10 \times (p_s/100)^{0.2}$, so $S'(100) = 12$. Then
$F_p = D' - S' = -17$ and $F_\tau = S' = 12$.

\textbf{4.} $\dd p/\dd\tau = -F_\tau/F_p = 12/17 = 0.7059$ kests of price per
kest of levy; $\dd Q/\dd\tau = D' \times 0.7059 = -3.53$ thousand tonnes per
kest.

\textbf{5.} $0.7059 - 1 = -0.2941$: each kest of levy lowers the price
producers keep by 0.29 kests.

\textbf{6.} Both lines give 1000 at $p = 100$ and have the slopes $-5$ and
$12$ of $D$ and $S$ there: they are the tangents. $1500 - 5p = 12p - 12\tau -
200$ gives $17p = 1700 + 12\tau$, $p = 100 + (12/17)\tau$.

\textbf{7.} $\varepsilon_d = -0.5$ and $\varepsilon_s = 1.2$, the exponents.

\textbf{8.} $1.2/(1.2 + 0.5) = 1.2/1.7 = 0.7059$, as in question 4.

\textbf{9.} $1 - 0.7059 = 0.2941$, about $29\%$.

\textbf{10.} More elastic supply (imports of cement from Brenland, idle
kilns) or less elastic demand.

\textbf{11.} $1 + \varepsilon_d = 0.5$.

\textbf{12.} The price rises by $0.7059 \times 10 = 7.06$ kests, $7.06\%$; spending
rises by $0.5 \times 7.06 = 3.53\%$.

\textbf{13.} The first-order condition $p_s = C'(Q) = 100(Q/1000)^{1/1.2}$ gives
$Q = 1000(p_s/100)^{1.2} = S(p_s)$. By part (1) of
\cref{thm:b1:comparative-statics:envelope}, $\Pi'(p_s) = \partial(p_sQ - C(Q))/\partial p_s = Q$
at the optimum, that is $S(p_s)$.

\textbf{14.} $\Delta p_s \approx -0.2941 \times 10 = -2.94$ kests; the loss of profit
is about $1000 \times 2.94 = 2940$ thousand kests.

\textbf{15.} Buyers: $1000 \times 7.06 = 7060$ thousand kests. The sum,
$10000$, equals the levy on 1000 thousand tonnes at 10 kests: to first
order the levy's burden is split between the two sides and equals its
revenue.

\textbf{16.} Smaller: $\Pi'' = S' > 0$, so $\Pi$ is convex; as $p_s$ falls,
producers sell less, and each further kest of price is lost on fewer
tonnes. The first-order loss counts every kest on the initial 1000.

\textbf{17.} $p \approx 107.06$, $p - \tau \approx 97.06$, and
$\Delta\ln Q \approx -0.5 \times 7.06\% = -3.53\%$, so $Q \approx 964.7$.

\textbf{18.} $D(107.16) = 1000 \times 1.0716^{-0.5} = 966.0$ and
$S(97.16) = 1000 \times 0.9716^{1.2} = 966.0$: the market clears. $Q = 966.0$
and $p - \tau = 97.16$.

\textbf{19.} The first-order price is $0.10$ kest too low and the first-order
quantity $1.3$ thousand tonnes too low.

\textbf{20.} $10 \times 966.0 = 9660$ thousand kests.

\textbf{21.} Before, $100 \times 1000 = 100000$; after, $107.16 \times 966.0 =
103517$, a rise of $3.52\%$, against $3.53\%$ predicted.

\textbf{22.} $\int_{97.16}^{100} 1000(s/100)^{1.2}\,\dd s = \frac{1000 \times 100}{2.2}
\bigl[1 - 0.9716^{2.2}\bigr] = 45454.5 \times 0.06142 = 2792$ thousand kests,
$148$ less than the first-order 2940.

\textbf{23.} The first-order prediction exactly: the tangent lines have the
derivatives of $D$ and $S$ at $(100, 1000)$, so the linear model's response is
the first-order response, $p = 100 + 0.7059 \times 10 = 107.06$.

\textbf{24.} To first order, 71 kest cents of every kest of levy reach the
price of cement ($0.7059$); exactly, for a levy of 10 kests, $7.16/10$, that is
72 cents.
\end{solution}
